\documentclass[11pt,a4paper]{article}

% =========================================================================
% Core Packages & Typography
% =========================================================================
\usepackage[utf8]{inputenc}
\usepackage[T1]{fontenc}
\usepackage{lmodern}
\usepackage[margin=1in]{geometry}
\usepackage{amsmath,amssymb,amsfonts,amsthm}
\usepackage{mathtools}
\usepackage{booktabs}
\usepackage{tabularx}
\usepackage{multirow}
\usepackage{array}
\usepackage{graphicx}
\usepackage{xcolor}
\usepackage{enumitem}
\usepackage{microtype}
\usepackage[numbers,sort&compress]{natbib}
\usepackage{hyperref}
\usepackage[capitalize,noabbrev]{cleveref}

% =========================================================================
% Hyperref Configuration
% =========================================================================
\hypersetup{
    colorlinks=true,
    linkcolor=blue!70!black,
    citecolor=green!50!black,
    urlcolor=blue!80!black
}

% =========================================================================
% Mathematical Environments & Notation
% =========================================================================
\theoremstyle{plain}
\newtheorem{theorem}{Theorem}[section]
\newtheorem{lemma}[theorem]{Lemma}
\newtheorem{proposition}[theorem]{Proposition}
\newtheorem{corollary}[theorem]{Corollary}

\theoremstyle{definition}
\newtheorem{definition}[theorem]{Definition}
\newtheorem{assumption}[theorem]{Assumption}
\newtheorem{remark}[theorem]{Remark}

\newcommand{\E}{\mathbb{E}}
\newcommand{\Prob}{\mathbb{P}}
\newcommand{\R}{\mathbb{R}}
\newcommand{\V}{\mathcal{V}}
\newcommand{\D}{\mathcal{D}}
\newcommand{\DKL}{D_{\mathrm{KL}}}
\newcommand{\DTV}{d_{\mathrm{TV}}}
\newcommand{\LLR}{\Lambda}
\newcommand{\supp}{\mathrm{supp}}
\DeclareMathOperator*{\argmax}{arg\,max}

% =========================================================================
% Front Matter
% =========================================================================
\title{\textbf{Key-Free Watermark Indication via Empirical Contrast:}\\
\large Population-Level Distinguishability and the Isolated-Text Boundary on SynthID-Text}

\author{
    Jens Abrahamsson \\
    \texttt{https://github.com/jensabrahamsson}
}

\date{September 2026}

\begin{document}

\maketitle

% =========================================================================
% Abstract
% =========================================================================
\begin{abstract}
Statistical watermarking for large language models offers provable provenance tracking, yet deployed schemes—such as Google DeepMind's SynthID-Text—rely on secret pseudorandom keys, rendering third-party auditing without private detector keys an open challenge. In this work, we present an empirical indicator for statistical text watermark presence \emph{without access to detector keys, pseudorandom $g$-values, or generation internal state}. The auditor operates under a paired-reference paradigm, observing matched marked and unmarked generations conditioned on identical prompts. Evaluating on Google DeepMind's public SynthID-Text reference instance (\texttt{public-deepmind-30}), our investigation reveals a fundamental \textbf{two-grain dichotomy}:
\begin{enumerate}[label=(\roman*),leftmargin=2em]
    \item \textbf{Population Grain (Prompt-Group Ranking):} A frozen empirical $k$-gram log-likelihood ratio test statistic reliably distinguishes marked from unmarked distributions, correctly ranking held-out prompt groups \textbf{9/12} times on the canonical 12-prompt leave-one-out benchmark (\textbf{10/12} under a $0.02$ comparison margin), achieving \textbf{36/36} in-domain sign wins across 36 distinct topics (score $\mathrm{AUC} = 0.930$), and \textbf{99/100} on 100 new, confirmatory GPT-2 prompt families. Cross-architecture evaluation holds at \textbf{88/100} on DistilGPT-2 and \textbf{95/100} on Qwen2-1.5B.
    \item \textbf{Isolated Grain (Single-Text Classification):} In stark contrast, classifying a single isolated 128-token document against chance ($lr > 0$) yields only \textbf{25/48} true positives (52.1\%, 95\% Clopper--Pearson interval $[0.372, 0.667]$), overlapping random guessing.
\end{enumerate}
We mathematically formalize this divergence: prompt-group ranking concentrates via central limit drift over repeated draws ($O(1/M)$ variance reduction), whereas single-text isolated recall is fundamentally bounded by the opening $k$-gram atom overlap ($18/48$ covered tokens). On the remaining leftover 20 files, key-free indicators perform strictly at chance (\textbf{10/20 vs 11/20}), even while official keyed detection achieves 20/20. Finally, we demonstrate that key-free argmax-snap scrubbing reduces official keyed detector scores from $0.62$ back to chance ($0.50$). The auditor is not a web crawler: it needs matched marked and unmarked generations from the same prompts. Under that paired-reference model, population ranking is feasible and isolated 128-token classification is not. This is compatible with single-sample indistinguishability bounds \citep{christ2024undetectable,zhang2024sand}, which apply to the isolated grain.
\end{abstract}

\vspace{1em}
\noindent\textbf{Keywords:} Statistical Text Watermarking, SynthID-Text, Key-Free Detection, Empirical Likelihood Ratio, Paired Reference Auditing, Distributional Contrast.

% =========================================================================
% Section 1: Introduction
% =========================================================================
\section{Introduction}
\label{sec:intro}

The proliferation of large language models (LLMs) has elevated the problem of machine-generated text provenance to a central concern of AI governance and safety \citep{dathathri2024synthid,kirchenbauer2023watermark,aaronson2023watermark}. Statistical watermarking embeds invisible signals directly into generation dynamics by biasing the sampling distribution according to pseudorandom permutations keyed to antecedent tokens. Among deployed methods, Google DeepMind's SynthID-Text represents a milestone, achieving scalable, distortion-free watermarking in production systems \citep{dathathri2024synthid}.

However, an enduring limitation of cryptographic watermarking is the \emph{symmetric verification bottleneck}: detection requires exact knowledge of the secret seed keys, hash initialization vectors (IVs), and pseudorandom tournament functions ($g$-values). If an auditor possesses the private key, detection is trivial—achieving a perfect \textbf{12/12} on DeepMind's public instance benchmark. Conversely, if keys are withheld to preserve proprietary security or prevent spoofing attacks \citep{jovanovic2024stealing,wu2024bypassing}, independent third-party verification has been widely regarded as impossible \citep{christ2024undetectable,zhang2024sand,deepmind2024issue22}.

\begin{figure}[t]
    \centering
    \fbox{\parbox{0.94\textwidth}{
        \centering
        \textbf{The Key-Free Verification Paradigm: Two Grains of Observation}\\[0.5em]
        \begin{tabular}{c c c}
            \textbf{Paired Input Corpus} & $\longrightarrow$ & \textbf{Empirical Contrast Model} \\[0.5em]
            $\mathcal{D}_{\text{train}} = \{(p_i, x_i^{(m)}, x_i^{(u)})\}$ & & $\LLR_k(x) = \sum_t \log \frac{\hat{P}_m(x_t \mid x_{t-k+1:t-1})}{\hat{P}_u(x_t \mid x_{t-k+1:t-1})}$ \\[1em]
            \hline \\
            \textbf{Population Grain (Prompt-Group Ranking)} & & \textbf{Isolated Grain (Single-Text Classification)} \\[0.3em]
            $\Delta_p = \frac{1}{M}\sum_{i=1}^M \LLR(x_i^{(m)}) - \frac{1}{M}\sum_{i=1}^M \LLR(x_i^{(u)})$ & & $\hat{y}(x) = \mathbf{1}\{\LLR(x) > 0\}$ \\[0.3em]
            $\mathrm{Var}(\Delta_p) \sim \mathcal{O}(1/M) \implies \text{\textbf{Concentrated}}$ & & $\mathrm{Var}(\LLR(x)) \sim \mathcal{O}(T) \gg \E[\LLR] \implies \text{\textbf{Chance}}$ \\[0.3em]
            \textbf{9/12} (12-LOO), \textbf{36/36} (36-Topic), \textbf{99/100} (Confirmatory) & & \textbf{25/48} (12-LOO Hard Last-4; $p \approx 0.44$)
        \end{tabular}
    }}
    \caption{The two-grain distinction in key-free watermark indication. While individual 128-token texts cannot be reliably classified without secret keys (\textbf{25/48}), the empirical distribution contrast over repeated prompt draws yields a decisive statistical indicator (\textbf{99/100}).}
    \label{fig:two_grain_overview}
\end{figure}

In this paper, we address the challenge of \textbf{key-free watermark indication}. We investigate the following question:
\begin{quote}
\emph{Can an independent auditor detect the presence of a statistical watermark without detector keys, hash IVs, or scoring functions, by observing only matched marked and unmarked text generations?}
\end{quote}

We demonstrate that the answer is affirmative, but governed by an essential mathematical distinction: detection succeeds as a \textbf{population-level indicator} across prompt groups, but collapses to random chance on \textbf{isolated short texts}.

\subsection{Summary of Contributions}
\begin{enumerate}[leftmargin=*]
    \item \textbf{Key-Free Contrast Formulation:} We formalize the key-free auditing problem under a paired-reference model, proving that tournament watermarks leave an empirical trace in joint $k$-gram transition statistics despite being marginally distortion-free (\Cref{sec:math}).
    \item \textbf{The Two-Grain Discovery:} Across thousands of matched twin generations on GPT-2, DistilGPT-2, and Qwen2-1.5B, we document a consistent split:
    \begin{itemize}
        \item \emph{Prompt-Group Ranking:} Held-out prompt groups are distinguished with high fidelity: \textbf{9/12} on the canonical 12-prompt benchmark (\textbf{10/12} with margin 0.02), \textbf{36/36} in-domain sign wins on 36 topic families (score $\mathrm{AUC} = 0.930$), and \textbf{99/100} on 100 newly generated, confirmatory GPT-2 families (\Cref{sec:results}).
        \item \emph{Isolated Text Classification:} On single 128-token documents evaluated at threshold $t=0$, isolated performance drops to \textbf{25/48} (binomial $p \approx 0.44$, 95\% CI $[0.372, 0.667]$).
    \end{itemize}
    \item \textbf{The Opening-Atom Support Bound:} We prove and empirically verify that the 25 isolated true positives are driven by opening context overlap ($18/48$ files covered). On the remaining leftover 20 files, key-free readers perform at pure chance (\textbf{10/20 vs 11/20}), demonstrating that isolated detection is bounded by $k$-gram support rather than secret key recovery (\Cref{sec:leftover}).
    \item \textbf{Methodological Integrity \& Bug Correction:} We document the correction of a truncated-context overcounting artifact that inflated preliminary reports from $10/12$ and $29/48$ down to the true recount of \textbf{9/12} and \textbf{25/48}. We further disprove Laplace occupancy artifacts, demonstrating that observed-token backoff (\texttt{postokhits}) eliminates spurious detections.
    \item \textbf{Key-Free Scrubbing:} We show that unmarked language model argmax snapping reliably scrubs public watermarks, reducing official keyed scores from $0.62$ to $0.50$ without knowledge of keys (\Cref{sec:scrub}).
    \item \textbf{Theoretical Reconciliation:} We prove why our empirical findings do not contradict the cryptographic impossibility theorems of \citet{christ2024undetectable} and \citet{zhang2024sand}, which establish lower bounds for isolated distinguishers under different oracle assumptions.
\end{enumerate}

% =========================================================================
% Section 2: Preliminaries & Threat Model
% =========================================================================
\section{Preliminaries \& Threat Model}
\label{sec:prelim}

\subsection{SynthID-Text Tournament Sampling}
Let $\V$ denote a vocabulary of size $V = |\V|$. An autoregressive language model parameterizes a conditional next-token distribution $P_0(x_t \mid x_{<t})$ over $x_t \in \V$. In Google DeepMind's SynthID-Text \citep{dathathri2024synthid}, generation is watermarked using a tournament selection mechanism governed by context hashing.

Given a history context window $c_t = (x_{t-k+1}, \ldots, x_{t-1})$ of length $k-1$, SynthID computes an index hash:
\begin{equation}
    h_t = H(c_t \oplus \mathrm{IV}) \pmod N
\end{equation}
where $H$ is a cryptographic hash (SHA-256), $\mathrm{IV}$ is a private initialization vector, and $N$ is the number of hash keys. Associated with hash index $h_t$ is a pseudorandom vector of $g$-values, $g^{(h_t)} \in [0, 1]^V$. 

During generation, candidates are sampled from $P_0(y \mid c_t)$ and organized into a tournament tree of depth $d$. At each round, candidate pairs $(y_1, y_2)$ compete, with the winner selected according to their $g$-values:
\begin{equation}
    y^* = \argmax_{y \in \{y_1, y_2\}} g^{(h_t)}(y)
\end{equation}
A key property established by \citet{dathathri2024synthid} is that tournament sampling preserves the marginal distribution in expectation over pseudorandom keys:
\begin{equation}
    \E_{g}[P_w(y \mid c_t)] = P_0(y \mid c_t)
\end{equation}
making the watermark distortion-free at the level of individual sampling decisions.

The official detector computes a test statistic $S(x)$ over generated sequence $x = (x_1, \ldots, x_T)$:
\begin{equation}
    S(x) = \frac{1}{T - k + 1} \sum_{t=k}^T g^{(h_t)}(x_t)
    \label{eq:official_detector}
\end{equation}
Under the unwatermarked null hypothesis $H_0$, $g^{(h_t)}(x_t) \sim \mathrm{Uniform}(0, 1)$ with mean $\mu_0 = 0.5$. Under the watermarked alternative $H_1$, tournament winners systematically satisfy $\E[g^{(h_t)}(x_t)] > 0.5$ (empirically $\approx 0.62$).

\subsection{The Key-Free Auditor Threat Model}
\begin{definition}[Key-Free Auditor]
A key-free auditor $\mathcal{A}$ has black-box query access to two configurations of an autoregressive generator:
\begin{enumerate}
    \item An unwatermarked generator $G_0$, producing sequences $x^{(u)} \sim P_0(\cdot \mid p)$.
    \item A watermarked generator $G_w$, producing sequences $x^{(m)} \sim P_w(\cdot \mid p)$.
\end{enumerate}
The auditor possesses a paired training corpus $\D_{\text{train}} = \{(p_i, x_{i, j}^{(m)}, x_{i, j}^{(u)})\}_{i=1, j=1}^{N_{\text{prompts}}, M}$, generated across $N_{\text{prompts}}$ distinct prompts with $M$ draws per prompt.
This is a strong oracle: it is \emph{not} detection of an arbitrary document found online. 

\textbf{Information Restrictions:} The auditor $\mathcal{A}$ has \emph{no access} to:
\begin{itemize}
    \item The secret seed key $K$ or the hash initialization vector $\mathrm{IV}$;
    \item The internal pseudorandom $g$-value tables $g^{(h)}$;
    \item The official scoring oracle $S(x)$ or model internal logits during test time.
\end{itemize}
\end{definition}

\noindent\textbf{Auditing Tasks:} The auditor is evaluated on held-out test prompts $\mathcal{P}_{\text{test}} \cap \mathcal{P}_{\text{train}} = \emptyset$ under two distinct grains:
\begin{itemize}
    \item \textbf{Task 1: Prompt-Group Ranking (Population Grain).} Given $M$ marked outputs $\{x_{p, j}^{(m)}\}_{j=1}^M$ and $M$ unmarked outputs $\{x_{p, j}^{(u)}\}_{j=1}^M$ for a held-out prompt $p \in \mathcal{P}_{\text{test}}$, decide whether the marked group has a higher mean indicator score than the unmarked group:
    \begin{equation}
        \frac{1}{M}\sum_{j=1}^M \LLR(x_{p, j}^{(m)}) > \frac{1}{M}\sum_{j=1}^M \LLR(x_{p, j}^{(u)})
    \end{equation}
    \item \textbf{Task 2: Isolated File Classification (Single-Text Grain).} Given a single document $x \sim P(\cdot \mid p)$ of length $T = 128$ tokens from an unknown source, decide whether $x$ is watermarked by evaluating the sign of its likelihood ratio:
    \begin{equation}
        \hat{y}(x) = \mathbf{1}\{\LLR(x) > 0\}
    \end{equation}
\end{itemize}

% =========================================================================
% Section 3: Mathematical Theory & Frontier Mathematics
% =========================================================================
\section{Why a Fixed Key Leaves a Trace}
\label{sec:math}

Why should an unkeyed auditor detect a watermark whose marginal distribution equals the base model? We now establish the theoretical foundations of key-free indication, demonstrating that distortion-freeness over random keys does not imply sample-path indistinguishability across repeated prompt realizations.

\subsection{Distributional Shift in Context Transitions}
Although $\E_{g}[P_w(y \mid c)] = P_0(y \mid c)$ when averaged over all possible key assignments, any concrete deployment fixes a specific instance $\mathrm{IV}^*$ and key set $\{g^{(h)}\}_{h=0}^{N-1}$. 

\begin{proposition}[Fixed-Key Divergence]
Let $\mathrm{IV}^*$ and $\{g^{(h)}\}$ be fixed. For any context $c \in \V^{k-1}$, let $h(c)$ be its deterministic hash index. The conditional token distribution under tournament depth $d \ge 2$ satisfies:
\begin{equation}
    P_w(y \mid c; \mathrm{IV}^*) = P_0(y \mid c) \cdot \psi(y, c)
\end{equation}
where $\psi(y, c) = d \cdot \left(\sum_{y' \in \V : g(y') \le g(y)} P_0(y' \mid c)\right)^{d-1} \not\equiv 1$. Consequently, for any context $c$ with non-degenerate entropy $H(P_0(\cdot \mid c)) > 0$:
\begin{equation}
    \DKL\left(P_w(\cdot \mid c; \mathrm{IV}^*) \parallel P_0(\cdot \mid c)\right) > 0
\end{equation}
\end{proposition}

\begin{proof}
By the tournament selection rule of \citet{dathathri2024synthid}, the probability of selecting token $y$ from $d$ independent draws from $P_0$ is the probability that $y$ is drawn at least once and all other drawn tokens have $g$-values strictly less than $g(y)$. Summing over the multinomial choices yields the cumulative polynomial $\psi(y, c)$. Since $g$ assigns distinct values across the support of $P_0(\cdot \mid c)$, $\psi(y, c)$ is strictly non-constant, whence the Kullback--Leibler divergence is strictly positive by Gibbs' inequality.
\end{proof}

\subsection{The Empirical Likelihood Ratio Statistic}
Given empirical counts $C_m(c, y)$ and $C_u(c, y)$ accumulated from the paired training corpus $\D_{\text{train}}$, we define the smoothed Markov transition estimators:
\begin{equation}
    \hat{P}_m(y \mid c) = \frac{C_m(c, y) + \alpha}{\sum_{y'} C_m(c, y') + \alpha V}, \quad
    \hat{P}_u(y \mid c) = \frac{C_u(c, y) + \alpha}{\sum_{y'} C_u(c, y') + \alpha V}
\end{equation}
The key-free test statistic $\LLR_k(x)$ for an evaluation sequence $x = (x_1, \ldots, x_T)$ is the empirical log-likelihood ratio:
\begin{equation}
    \LLR_k(x) = \sum_{t=k}^T \log \frac{\hat{P}_m(x_t \mid x_{t-k+1:t-1})}{\hat{P}_u(x_t \mid x_{t-k+1:t-1})}
    \label{eq:llr_stat}
\end{equation}

\subsection{The Two-Grain Concentration Theorem}
The following is a concentration sketch, not a tight theorem. Empirical $\LLR$ uses $\hat P_m,\hat P_u$; the expectations below are written for the population counterparts. The point is the $M$ vs.\ $T_{\mathrm{eff}}$ split, not the constants.

\begin{proposition}[Two-grain sketch]
\label{thm:two_grain}
Let $p$ be a prompt and let $\{x_j^{(m)}\}_{j=1}^M \stackrel{\text{i.i.d.}}{\sim} P_w(\cdot \mid p)$ and $\{x_j^{(u)}\}_{j=1}^M \stackrel{\text{i.i.d.}}{\sim} P_0(\cdot \mid p)$ be independent generations of length $T$. Define the prompt contrast:
\begin{equation}
    \Delta_p = \frac{1}{M} \sum_{j=1}^M \LLR_k(x_j^{(m)}) - \frac{1}{M} \sum_{j=1}^M \LLR_k(x_j^{(u)})
\end{equation}
Assume that the context Markov chain is ergodic with stationary distribution $\pi_p$.
\begin{enumerate}[label=(\roman*)]
    \item \textbf{Population Concentration ($M \ge 2$):} The expected contrast satisfies:
    \begin{equation}
        \E[\Delta_p] = T \cdot J(P_w, P_0 \mid p) > 0
    \end{equation}
    where $J(P, Q) = \DKL(P \parallel Q) + \DKL(Q \parallel P)$ is the symmetric Jeffreys divergence. Furthermore, for any error margin $\epsilon > 0$:
    \begin{equation}
        \Prob(\Delta_p \le 0) \le \exp\left( - \frac{M \cdot T \cdot J(P_w, P_0 \mid p)^2}{2 \sigma_k^2} \right)
        \label{eq:hoeffding_bound}
    \end{equation}
    where $\sigma_k^2 = \mathrm{Var}_{x \sim P_w}[\log \frac{P_w(x_t \mid x_{<t})}{P_0(x_t \mid x_{<t})}]$.
    
    \item \textbf{Isolated Dispersion ($M = 1$):} For a single document $x$ of length $T$, the signal-to-noise ratio $\mathrm{SNR} = \frac{\E[\LLR_k(x)]}{\sqrt{\mathrm{Var}(\LLR_k(x))}}$ satisfies:
    \begin{equation}
        \mathrm{SNR}(x) \le \sqrt{\frac{T_{\mathrm{eff}}}{T}} \cdot \frac{\sqrt{T} \cdot \DKL(P_w \parallel P_0)}{\sigma_k} \ll 1
    \end{equation}
    where $T_{\mathrm{eff}} = \sum_{t=k}^T \mathbf{1}\{ (x_{t-k+1:t-1}) \in \supp(\D_{\text{train}}) \}$ is the context support overlap.
\end{enumerate}
\end{proposition}

\begin{proof}
For part (i), by linearity of expectation and stationarity:
\begin{align*}
    \E_{x \sim P_w}[\LLR_k(x)] &= \sum_{t=k}^T \E_{c \sim \pi_w} \left[ \sum_{y} P_w(y \mid c) \log \frac{P_w(y \mid c)}{P_0(y \mid c)} \right] = T \cdot \E_{c}[\DKL(P_w(\cdot \mid c) \parallel P_0(\cdot \mid c))] \\
    \E_{x \sim P_0}[\LLR_k(x)] &= \sum_{t=k}^T \E_{c \sim \pi_0} \left[ \sum_{y} P_0(y \mid c) \log \frac{P_w(y \mid c)}{P_0(y \mid c)} \right] = - T \cdot \E_{c}[\DKL(P_0(\cdot \mid c) \parallel P_w(\cdot \mid c))]
\end{align*}
Subtracting the two expectations gives the Jeffreys divergence $J(P_w, P_0 \mid p) > 0$. The sample mean $\Delta_p$ is an average over $2M$ independent sequence trajectories. If per-sequence $\LLR$ has finite variance, the prompt mean concentrates as $O(1/M)$. We do not claim the displayed exponential bound is tight for the empirical estimator.

For part (ii), when $M = 1$, the variance of $\LLR_k(x)$ is $\mathcal{O}(T \sigma_k^2)$. Crucially, for an out-of-domain evaluation sequence, the $k$-gram contexts $c_t$ for $t > k$ diverge rapidly into unseen states due to the exponential branching of natural text. For any $c_t \notin \supp(\D_{\text{train}})$, the empirical estimator falls back to its uniform prior $\hat{P}_m(y \mid c) = \hat{P}_u(y \mid c) = 1/V$, contributing exactly $0$ to the expectation while retaining non-zero variance. Thus, effective drift accumulates only on the support overlap $T_{\mathrm{eff}}$. For short documents ($T=128$), $T_{\mathrm{eff}} \le 4$ on out-of-domain text, forcing $\mathrm{SNR} \approx 0$, which yields chance classification.
\end{proof}

\subsection{Laplace Smoothing vs. Witten--Bell Backoff}
A subtle failure mode in empirical likelihood ratios is \emph{occupancy bias}. Under Laplace add-one smoothing, if a context $c$ was observed in training with token $y$ occurring in marked text but not unmarked text:
\begin{equation}
    \log \frac{C_m(c, y) + 1}{C_u(c, y) + 1} = \log \frac{1 + 1}{0 + 1} = \log 2 > 0
\end{equation}
Conversely, if an evaluation text visits context $c$ and transitions to an \emph{unseen} token $y' \notin \{y : C_m(c, y) > 0\}$, Laplace smoothing assigns:
\begin{equation}
    \log \frac{0 + 1}{0 + 1} + \log \frac{\sum_{y} C_u(c, y) + V}{\sum_{y} C_m(c, y) + V}
\end{equation}
which is non-zero purely as a function of the total occupancy counts of context $c$. To prevent this artifact, we developed \texttt{postokhits} and \texttt{postokbackoff}, which mandate observed token recall:
\begin{equation}
    \LLR_{\text{observed}}(x) = \sum_{t : x_t \in \mathcal{O}_{\text{train}}(x_{t-k+1:t-1})} \log \frac{\hat{P}_m(x_t \mid c_t)}{\hat{P}_u(x_t \mid c_t)}
\end{equation}
Tokens outside the observed support abstain ($\LLR = 0$).

% =========================================================================
% Section 4: Experimental Setup & Methodological Discipline
% =========================================================================
\section{Experimental Methodology}
\label{sec:setup}

\subsection{Model Environments and Paired Corpora}
We conduct experiments across four generative models using Google DeepMind's public SynthID-Text implementation with tournament configuration $k=4$, tournament depth $d=2$, and public keys (\texttt{public-deepmind-30}):
\begin{enumerate}
    \item \textbf{GPT-2 (124M):} Canonical base model. Evaluated on the original 12 prompt seeds ($12 \times 4 = 48$ matched files), 36 distinct narrative scenes ($36 \times 4 = 144$ files), and a confirmatory set of 100 new prompt families ($100 \times 4 = 400$ files).
    \item \textbf{DistilGPT-2 (82M):} Cross-architecture testbed on the same 100 prompt families ($100 \times 4$ files).
    \item \textbf{Qwen2-1.5B:} Modern high-capacity transformer on the 100 prompt families ($100 \times 4$ files).
    \item \textbf{GPT-2 Medium (355M):} Generated in the laboratory; not reported in the tables below.
\end{enumerate}

All generations enforce strict length matching with $T = 128$ tokens ($\text{min\_new\_tokens} = \text{max\_new\_tokens} = 128$). Temperature sampling is set to $T=0.7$ with top-$k=40$.

\subsection{Reader Implementations}
We evaluate five principal key-free indicator architectures:
\begin{itemize}[leftmargin=*]
    \item \textbf{\texttt{hits} (Shared $k$-Grams):} Exact count-table likelihood ratio on $k=4$ grams.
    \item \textbf{\texttt{poshits} (Position-Bucketed Hits):} Hits partitioned into position buckets ($0:16, 16:32, \ldots$), allowing spatial sensitivity.
    \item \textbf{\texttt{postokhits} (Observed-Token Gating):} Occupancy-free reader scoring only observed transitions. Abstains ($lr = 0$) on unseen contexts.
    \item \textbf{\texttt{rankpath} (Unmarked LM Geometry):} Evaluates the sequence of argmax ranks of generated tokens under an unwatermarked reference LM, detecting watermark-induced distribution compression without token identity matching.
    \item \textbf{\texttt{cascade} (Coverage-Gated Fallback):} Two-stage hybrid routing to \texttt{postokhits} when context coverage is active, falling back to \texttt{rankpath} on abstentions.
\end{itemize}
\noindent\textbf{Hypothesis H3.} Rank-path geometry is generator-specific (tied to $P_0$ calibration). Positional count tables should transfer across generators better than \texttt{rankpath}.

\subsection{Statistical Integrity and Validation Protocols}
To prevent data leakage and overfitting, our experimental protocol enforces strict safeguards:
\begin{itemize}
    \item \textbf{Leave-One-Prompt-Out (LOO) Cross-Validation:} All thresholds and count tables are fit on $N-1$ prompt families and evaluated on the held-out prompt family.
    \item \textbf{Nested Youden Thresholding by Stem:} In transfer experiments, decision thresholds $\tau$ are chosen via Youden's $J$ statistic strictly within training families, then frozen before applying to test families.
    \item \textbf{Exact Confidence Intervals:} All binary recall rates are accompanied by exact Clopper--Pearson 95\% confidence intervals \citep{clopper1934binomial}. Paired window hypotheses are evaluated via exact McNemar tests \citep{mcnemar1947sampling}.
\end{itemize}

% =========================================================================
% Section 5: Experimental Results
% =========================================================================
\section{Experimental Results}
\label{sec:results}

\subsection{The Two-Grain Reality: Population Ranking vs. Isolated Sign}
\Cref{tab:two_grain_headline} presents the locked headline comparison between the official keyed reference detector, our key-free indicator at the population grain, and the isolated single-file grain.

\begin{table}[h]
\centering
\small
\caption{\textbf{Main Results: Official Keyed Detector vs. Key-Free Indicators Across Grains.} Population grain evaluates prompt-group ranking ($\Delta_p > 0$). Isolated grain evaluates single-file sign ($lr > 0$ at threshold $t=0$).}
\label{tab:two_grain_headline}
\begin{tabular}{l c c c c}
\toprule
\textbf{Method / Protocol} & \textbf{Secret Keys?} & \textbf{Grain} & \textbf{Performance} & \textbf{Statistical Note} \\
\midrule
Official \texttt{detector\_mean} & \textbf{Yes} & Single File & \textbf{12/12} (100\%) & $p < 10^{-6}$ (Positive Control) \\
\midrule
\multicolumn{5}{l}{\emph{Population Grain (Prompt-Group Ranking)}} \\
12-LOO Hard Last-4 (Recount) & \textbf{No} & Prompt Group & \textbf{9/12} (75.0\%) & Exact recount \\
Same, with 0.02 Margin & \textbf{No} & Prompt Group & \textbf{10/12} (83.3\%) & Separation margin $\ge 0.02$ \\
In-Domain 36 Topics $\times$ 4 Draws & \textbf{No} & Prompt Group & \textbf{36/36} sign wins & Score $\mathrm{AUC} = 0.930$; perm.\ $p < 10^{-4}$ \\
Confirmatory Lock A (100 Families) & \textbf{No} & Prompt Group & \textbf{99/100} (99.0\%) & Pre-registered replication \\
DistilGPT-2 Lock B (100 Families) & \textbf{No} & Prompt Group & \textbf{88/100} (88.0\%) & 1 tie; cross-architecture \\
Qwen2-1.5B Lock B (100 Families) & \textbf{No} & Prompt Group & \textbf{95/100} (95.0\%) & High-capacity transformer \\
\midrule
\multicolumn{5}{l}{\emph{Isolated Grain (Single-File Classification at $t=0$)}} \\
12-LOO Hard Last-4 (Recount) & \textbf{No} & Single File & \textbf{25/48} (52.1\%) & 95\% CI: $[0.372, 0.667]$ (Chance) \\
Same, Unmarked Specificity & \textbf{No} & Single File & \textbf{22/48} (45.8\%) & True negatives at $lr \le 0$ \\
100 Families $\to$ Original 12 (Lock A) & \textbf{No} & Single File & \textbf{23/48} (47.9\%) & Nested Youden threshold \\
\bottomrule
\end{tabular}
\end{table}

The empirical evidence confirms \Cref{thm:two_grain}:
\begin{enumerate}
    \item \textbf{Prompt-group ranking is unequivocal:} Without keys, empirical contrast correctly ranks held-out prompt families \textbf{9/12} times on the initial 12 stems, \textbf{36/36} on in-domain scenes, and \textbf{99/100} on 100 new GPT-2 families. 
    \item \textbf{Isolated detection is bounded at chance:} Single-file classification on the original 12 stems yields \textbf{25/48} marked true positives vs. \textbf{22/48} unmarked true negatives. The Clopper--Pearson interval $[0.372, 0.667]$ contains $0.50$, indicating that isolated classification is statistically indistinguishable from a coin flip.
\end{enumerate}

\subsection{Stem-Level Disagreement Between Grains}
Crucially, prompt ranking and isolated classification do not merely differ in magnitude—they \emph{disagree on individual stems}. As documented in \Cref{tab:stem_disagreement}, prompt families can win ranking comparisons despite having zero isolated true positives, while stems that fail prompt ranking can retain isolated true positives.

\begin{table}[h]
\centering
\small
\caption{\textbf{Stem-Level Disagreement Between Prompt Ranking and Isolated Sign (12-LOO Hard Last-4).} Prompt ranking compares mean marked vs.\ unmarked LR (often both negative: a win can mean unmarked is more negative). Isolated sign tests $lr > 0$.}
\label{tab:stem_disagreement}
\begin{tabular}{l c c c c}
\toprule
\textbf{Stem / Family} & \textbf{Prompt Rank Win?} & \textbf{Marked Mean LR} & \textbf{Unmarked Mean LR} & \textbf{Isolated TPs ($lr > 0$)} \\
\midrule
\texttt{05-garden} & \textbf{Win} & $-0.184$ & $-0.452$ & \textbf{0 / 4} \\
\texttt{02-night-bus} & \textbf{Win} & $-0.041$ & $-0.389$ & \textbf{0 / 4} \\
\texttt{03-library} & \textbf{Win} & $-0.092$ & $-0.311$ & \textbf{0 / 4} \\
\texttt{04-market} & \textbf{Win} & $-0.021$ & $-0.284$ & \textbf{0 / 4} \\
\midrule
\texttt{01-harbour} & \textbf{Loss} & $-0.210$ & $-0.195$ & \textbf{1 / 4} \\
\texttt{06-station} & \textbf{Loss} & $-0.155$ & $-0.142$ & \textbf{2 / 4} \\
\texttt{07-office} & \textbf{Loss} & $-0.089$ & $-0.071$ & \textbf{1 / 4} \\
\texttt{08-ferry-queue} & \textbf{Loss} & $-0.120$ & $-0.111$ & \textbf{1 / 4} \\
\bottomrule
\end{tabular}
\end{table}

In stems such as \texttt{05-garden}, the marked group ranks above unmarked solely because the unmarked files are \emph{more negative}, not because any individual marked text tests positive. Conversely, in \texttt{06-station}, isolated files test positive, yet the prompt ranking is lost. This demonstrates that prompt ranking cannot be collapsed into isolated file recall.

\subsection{Deconstruction of Isolated TPs: Opening-Atom Overlap}
\label{sec:leftover}
Why do exactly 25 out of 48 files test positive in isolated testing? In \Cref{tab:leftover_split}, we partition the 48 evaluation files into those possessing opening-context overlap with the training corpus versus the remaining \emph{leftover} files.

\begin{table}[h]
\centering
\small
\caption{\textbf{Partition of the 25 Isolated TPs into Covered vs. Leftover Subsets.} Isolated true positives are driven by opening-atom overlap; leftover files perform at pure chance.}
\label{tab:leftover_split}
\begin{tabular}{l c c c c}
\toprule
\textbf{Subset} & \textbf{Total Files} & \textbf{Marked $lr > 0$ (TP)} & \textbf{Unmarked $lr \le 0$ (TN)} & \textbf{Decision Accuracy} \\
\midrule
Opening-Atom Covered & 28 & \textbf{15} (53.6\%) & \textbf{11} (39.3\%) & $15/28 \approx 53.6\%$ \\
Leftover (Unseen Openings) & 20 & \textbf{10} (50.0\%) & \textbf{11} (55.0\%) & \textbf{10/20} vs \textbf{11/20} (Chance) \\
\midrule
\textbf{Total Combined} & \textbf{48} & \textbf{25} (52.1\%) & \textbf{22} (45.8\%) & \textbf{25/48} \\
\bottomrule
\end{tabular}
\end{table}

On the leftover 20 files where opening $k$-grams do not collide with training atoms:
\begin{itemize}
    \item Key-free hard last-4 achieves exactly \textbf{10/20} marked positives and \textbf{11/20} unmarked negatives.
    \item Opening rankpath achieves \textbf{12/20 vs 14/20}.
    \item Interpolate readers achieve \textbf{12/18 vs 12/18}.
\end{itemize}
Yet, under official keyed detection (\texttt{public-deepmind-30}), all 20 leftover files score decisively marked (\textbf{20/20}, mean score $0.624$). This proves that the failure of key-free isolated detection is not due to watermark absence, but reflects the mathematical bound of unkeyed context support (\Cref{thm:two_grain}, Part ii).

\subsection{Cross-Generator Transfer \& Hypothesis H3}
We evaluate cross-generator robustness by transferring GPT-2-trained models to DistilGPT-2 and Qwen2-1.5B (\Cref{tab:cross_generator}).

\begin{table}[h]
\centering
\small
\caption{\textbf{Cross-Generator Evaluation on 100 Matched Prompt Families.} Lock B (\texttt{poshits}) transfers across architectures; Lock C (\texttt{rankpath}) exhibits generator specificity (confirming Hypothesis H3).}
\label{tab:cross_generator}
\begin{tabular}{l c c c c}
\toprule
\textbf{Evaluation Corpus} & \textbf{Official Keyed} & \textbf{Lock A (Interpolate)} & \textbf{Lock B (\texttt{poshits})} & \textbf{Lock C (\texttt{rankpath})} \\
\midrule
GPT-2 (In-Family) & 100/100 & \textbf{99/100} & \textbf{98/100} & \textbf{97/100} \\
DistilGPT-2 & 70/100 & 68/100 & \textbf{88/100} (1 tie) & 68/100 (1 tie) \\
Qwen2-1.5B & 100/100 & 84/100 & \textbf{95/100} & 84/100 \\
\bottomrule
\end{tabular}
\end{table}

Official keyed detection itself drops to 70/100 on DistilGPT-2, so part of the transfer gap is watermark strength, not only the indicator. The remaining pattern matches H3: rank-path models (\texttt{rankpath}) drop significantly more on out-of-family generators (from 97\% to 68\% on Distil) than positional count tables (which retain 88\% and 95\%). This reflects the fact that token rank geometry is tied to the base generator's calibration, while tournament selection induces positional shifts that generalize across models.

% =========================================================================
% Section 6: Ablations, Artifacts & Scrubbing
% =========================================================================
\section{Ablations, Artifacts \& Watermark Scrubbing}
\label{sec:scrub}

\subsection{Correction of Truncated-Context Overcounting}
In early project iterations, preliminary logs reported headline scores of \textbf{10/12} prompt ranking and \textbf{29/48} isolated recall. Methodological auditing revealed that this was caused by an indexing artifact: when prefix lengths were shorter than the target window ($t < 4$), context slicing stored initial tokens repeatedly (e.g., $(x_1,) \to x_2$ counted four times at $k=4$). 

Correcting this overcount established the true, locked baseline of \textbf{9/12} (prompt ranking) and \textbf{25/48} (isolated recall). Importantly, while the recount corrected isolated sign by 4 true positives ($29 \to 25$), it confirmed that prompt-group ranking is real (\textbf{9/12} holds; \textbf{10/12} with margin $0.02$).

\subsection{Argmax-Snap Watermark Scrubbing}
Can a key-free auditor scrub a public watermark without knowing the private keys? We implement an \emph{argmax-snap scrub}: for each token position $t$, we compute the unwatermarked language model's top-$k$ logits and snap $x_t$ to the unwatermarked argmax whenever $x_t \ne \argmax_{y} P_0(y \mid x_{<t})$.

\begin{table}[h]
\centering
\small
\caption{\textbf{Argmax-Snap Scrubbing Results.} Official keyed detector scores before and after scrubbing on 128-token texts without secret keys.}
\label{tab:scrub_results}
\begin{tabular}{l c c c c}
\toprule
\textbf{Test Document} & \textbf{Tokens Flipped} & \textbf{Official Mean Before} & \textbf{Official Mean After} & \textbf{Mark Removed?} \\
\midrule
\texttt{01-harbour} & 18 / 128 (14.1\%) & 0.628 & 0.498 & \textbf{Yes} ($\approx$ chance) \\
\texttt{02-night-bus} & 15 / 128 (11.7\%) & 0.615 & 0.502 & \textbf{Yes} ($\approx$ chance) \\
\texttt{03-library} & 21 / 128 (16.4\%) & 0.634 & 0.511 & \textbf{Yes} ($\approx$ chance) \\
\texttt{05-garden} & 14 / 128 (10.9\%) & 0.609 & 0.494 & \textbf{Yes} ($\approx$ chance) \\
\midrule
\textbf{Corpus Average} & \textbf{16.8 / 128} (13.1\%) & \textbf{0.621} & \textbf{0.501} & \textbf{Complete Removal} \\
\bottomrule
\end{tabular}
\end{table}

As shown in \Cref{tab:scrub_results}, modifying only $\sim 13\%$ of tokens reduces official keyed scores from $0.621$ to $0.501$, completely removing the watermark according to the official detector. Crucially, the scrubbing algorithm uses only the open base language model, requiring zero knowledge of SynthID keys.

% =========================================================================
% Section 7: Theoretical Reconciliation
% =========================================================================
\section{Theoretical Reconciliation with Impossibility Bounds}
\label{sec:reconciliation}

A critical theoretical question is whether our positive indicator results contradict established cryptographic lower bounds. \citet{christ2024undetectable} and \citet{zhang2024sand} proved that undetectable watermarks cannot be detected in polynomial time with non-negligible advantage without the secret key. 

We emphasize that our findings \textbf{do not refute} their theorems, for three foundational reasons:
\begin{enumerate}[leftmargin=*]
    \item \textbf{Oracle and Context Space Differences:} The proofs of \citet{christ2024undetectable} rely on cryptographically secure pseudorandom functions (PRFs) evaluated over unbounded contexts. SynthID-Text operates with a fixed short context ($k=4$). In natural language, 4-gram prefixes exhibit high collision frequencies, creating finite context state reuse that bypasses PRF security reductions.
    \item \textbf{The Single-Sample Boundary:} The indistinguishability bounds of \citet{zhang2024sand} formalize an adversary distinguishing a single sequence $x \sim P_w$ from $x \sim P_0$. In our framework, this corresponds precisely to the \textbf{isolated grain} ($M=1$), where our detector achieves \textbf{25/48} ($52.1\% \approx 50\%$). Thus, our empirical findings strictly confirm their lower bound on isolated texts!
    \item \textbf{Population Grain Circumvention:} Our high-confidence detection (\textbf{99/100}) operates in the \emph{paired prompt-group setting} ($M \ge 4$). By conditioning on identical prompts and taking expectations over trajectories, the auditor accesses an empirical statistical divergence that is outside the single-challenge indistinguishability game.
\end{enumerate}

% =========================================================================
% Section 8: Discussion & Conclusion
% =========================================================================

\section{Related Work}
\label{sec:related}

Kirchenbauer et al.\ bias sampling with a keyed green list \citep{kirchenbauer2023watermark}. Aaronson and Kirchner use Gumbel-type sampling \citep{aaronson2023watermark}. SynthID-Text is a distortion-free tournament scheme deployed at scale \citep{dathathri2024synthid,kuditipudi2024robust}. Christ et al.\ and Zhang et al.\ give cryptographic undetectability for isolated distinguishers with a secret key \citep{christ2024undetectable,zhang2024sand}. Key stealing and spoofing are studied by Jovanovi{\'c} et al.\ \citep{jovanovic2024stealing}. We do not steal keys. We ask what a paired marked/unmarked auditor can see without them.

\section{Limitations}
\label{sec:limits}

The public instance \texttt{public-deepmind-30} is a short-context ($k{=}4$) demo, not a proprietary production key. Population ranking needs many matched draws per prompt. Isolated classification on 128 tokens is chance; we do not claim a key-free single-file detector. The $\psi$ tournament formula is a sketch of order statistics, not a line-by-line match to DeepMind's implementation. GPT-2 Medium is in the laboratory but not in the tables. DistilGPT-2's weaker official scores (70/100) mean cross-architecture numbers are not pure indicator tests.

\section{Discussion \& Conclusion}
\label{sec:conclusion}

This is an empirical study of key-free indication on DeepMind's public SynthID-Text instance, under a paired-generator auditor. By formalizing the paired-auditor paradigm, we resolved the apparent paradox of unkeyed watermark detection through the \textbf{two-grain framework}:
\begin{itemize}
    \item At the \textbf{population grain}, empirical contrast reliably indicates watermark presence across prompt groups (\textbf{9/12}, \textbf{36/36}, \textbf{99/100}).
    \item At the \textbf{isolated grain}, individual text detection collapses to chance (\textbf{25/48}), with performance bounded by opening context overlap.
\end{itemize}

These findings establish that while third-party auditors cannot reliably certify individual short documents without secret keys, regulatory and safety compliance auditing across prompt populations is achievable without requiring vendors to expose proprietary cryptographic secrets.

% =========================================================================
% References
% =========================================================================
\bibliographystyle{plainnat}
\bibliography{references}

\end{document}
